\section{任务与数据集}

\subsection{三任务定义}

本文将轴承预测性维护拆分为三个互补任务，如表 \ref{tab:task-definitions} 所示。RUL 关注连续寿命趋势，HealthState 关注退化阶段，EarlyFault 关注是否已经进入预警状态。三个任务的共同点是都来自同一条轴承时间序列；区别在于输出形式、评价指标和工程含义不同。

\begin{table}[H]
\centering
\caption{三任务定义}
\label{tab:task-definitions}
\small
\begin{tabularx}{\textwidth}{p{2.5cm}p{3.2cm}p{2.0cm}p{2.8cm}Y}
\toprule
任务 & 输出 & 类型 & 主要指标 & 工程作用 \\
\midrule
RUL & \code{linear_rul_norm} & 回归 & RMSE & 预测当前样本距离寿命终点的归一化剩余比例。 \\
HealthState & \code{health_state_id} & 四分类 & WeightedF1 & 将退化过程划分为健康、轻微退化、中度退化和严重退化阶段。 \\
EarlyFault & \code{early_fault} & 二分类 & WeightedF1 & 判断样本是否进入早期异常或早期故障预警阶段。 \\
\bottomrule
\end{tabularx}
\end{table}


RUL 的目标值为 \code{linear_rul_norm}。该目标用 1 表示寿命早期，用 0 表示寿命终点，中间按时间顺序线性下降。HealthState 和 EarlyFault 不是数据集原始给出的人工故障类别，而是根据退化过程和 FPT 等规则构造的伪标签。因此，分类结果在论文中应解释为“退化阶段识别”和“早期预警能力”，不能等同于人工拆检得到的真实故障类型标签。

\subsection{数据集}

本文使用 XJTU-SY 与 PHM2012 两个轴承振动数据集。二者都包含轴承从正常运行到失效附近的振动信号，但实验语义不同。XJTU-SY 更强调多工况条件下的全寿命退化过程，PHM2012 更强调预测挑战中的学习集和测试集划分。表 \ref{tab:dataset-summary} 总结了两个数据集在本文中的使用方式。

\begin{table}[H]
\centering
\caption{数据集与实验划分语义}
\label{tab:dataset-summary}
\small
\begin{tabularx}{\textwidth}{p{2.6cm}p{4.0cm}p{3.8cm}Y}
\toprule
数据集 & 数据特点 & 划分方式 & 使用目的 \\
\midrule
XJTU-SY & 多工况、全寿命轴承振动数据，包含不同转速和载荷条件。 & 按轴承划分训练、验证和测试，避免同一轴承时间片泄漏。 & 检验跨轴承和多工况条件下的特征与模型泛化能力。 \\
PHM2012 & 预测挑战数据集，包含学习集和测试集，常用于轴承寿命预测。 & 优先遵循官方学习集和测试集语义，并保持轴承内部时间顺序。 & 检验方法在经典 PHM 预测任务上的表现。 \\
\bottomrule
\end{tabularx}
\end{table}


\subsection{划分原则}

为了避免信息泄漏，本文以轴承为基本划分单位。也就是说，同一轴承的连续时间片不会同时出现在训练、验证和测试数据中。对于序列样本，窗口长度设置为 8，只在同一轴承和同一数据子集内部构造，不跨轴承、不跨数据子集拼接。这一点对 RUL 特别重要，因为相邻时间片之间高度相关，如果随机打散，会造成模型提前看到测试轴承的退化轨迹。

在 XJTU-SY 上，实验重点是跨轴承和多工况下的泛化能力。在 PHM2012 上，实验优先保留官方学习集与测试集的语义。后续所有结果解释都遵循这一划分逻辑。
